\section{Patch contributions and the representation theorem}
\label{sec:representation}

\subsection{One-patch contributions}

Let \(P\) be a labeled patch and \(X^P\) its local active indicator.
If \(P\) begins at an outgoing interaction, let \(\alpha_P\) be its
hidden split-or-birth kind and put
\[
  \sigma_P=\sigma_{i(P)}^{\alpha_P}(S(P));
\]
for an incoming patch, put \(\sigma_P=1\).  The full contribution is
\begin{equation}
  C(P)
  =
  \E_P^{\mathrm{con}}\left[
    \sigma_P
    \exp\left(
      V_{i(P)}
      \int_{s(P)}^{e(P)}X_u^P\,\dd u
    \right)
  \right],
  \label{eq:full-contribution}
\end{equation}
where the expression is interpreted by an endpoint limit when
\(e(P)=\infty\).

If \(P\) is an end patch at time \(t\), its terminal contribution at
\(z\in[0,1]\) is
\begin{equation}
  C(z,P)
  =
  \E_P^{\mathrm{con}}\left[
    \sigma_P
    \exp\left(
      V_{i(P)}
      \int_{s(P)}^tX_u^P\,\dd u
    \right)
    z^{X_t^P}
  \right].
  \label{eq:end-contribution}
\end{equation}
These factors are computed explicitly in
\cref{app:patch-calculations}.  In particular, every end contribution
is affine in \(z\).

\subsection{Random-product representation}

Write \(\E_A\) for expectation over the successful-interaction
skeleton of the signed dual started from \((A,+)\).

\begin{theorem}[Patch representation]
\label{thm:patch-representation}
For every \(A\Subset\Lambda\), \(t\ge0\), and
\(\xi\in\{0,1\}^{\Lambda}\),
\begin{equation}
  P_t\chi_A(\xi)
  =
  \E_A\left[
    \prod_{P\in\cB_t}C(P)
    \prod_{P\in\cE_t}C(\xi(i(P)),P)
  \right].
  \label{eq:patch-representation}
\end{equation}
\end{theorem}

\begin{proof}
The patches partition the active site-time set up to time \(t\).
Every successful nonempty-target interaction contributes its sign to
exactly one outgoing-initial patch, and every end patch records the
dual state at its site at time \(t\).  Consequently the integrand in
the Feynman--Kac formula \eqref{eq:fk-duality} decomposes exactly as
\begin{align*}
  &\sigma_t
  \exp\left(\int_0^tV(A_u)\,\dd u\right)
  \chi_{A_t}(\xi)
  \\
  &\qquad=
  \prod_{P\in\cB_t}
  \left[
    \sigma_P
    \exp\left(
      V_{i(P)}
      \int_{s(P)}^{e(P)}X_u^P\,\dd u
    \right)
  \right]
  \\
  &\qquad\quad\times
  \prod_{P\in\cE_t}
  \left[
    \sigma_P
    \exp\left(
      V_{i(P)}
      \int_{s(P)}^tX_u^P\,\dd u
    \right)
    \xi(i(P))^{X_t^P}
  \right].
\end{align*}
Apply \cref{thm:patch-factorization} conditionally on \(\cG_t\) and
then take expectation.
\end{proof}

\subsection{General initial laws}

Integrating \eqref{eq:patch-representation} against a probability
measure \(\mu\) gives
\begin{equation}
  \mu(P_t\chi_A)
  =
  \E_A\left[
    \prod_{P\in\cB_t}C(P)\,
    \mu\left(
      \prod_{P\in\cE_t}C(\eta(i(P)),P)
    \right)
  \right].
  \label{eq:general-law-representation}
\end{equation}

Suppose that a profile \(\pstar\) has been chosen so that every end
factor admits the affine decomposition
\begin{equation}
  C(z,P)
  =
  C(p_{i(P)}^\star,P)
  +
  b(P)\bigl(z-p_{i(P)}^\star\bigr),
  \qquad
  C(p_{i(P)}^\star,P)\ge0,\quad b(P)\ge0.
  \label{eq:centered-end-factor}
\end{equation}
Since different end patches are based at different sites, expansion
of the product yields
\begin{align}
  &\mu\left(
    \prod_{P\in\cE_t}C(\eta(i(P)),P)
  \right)
  \notag\\
  &\quad=
  \sum_{\cQ\subseteq\cE_t}
  \mu\left(
    \chi_{\{i(P):P\in\cQ\}}^\star
  \right)
  \prod_{P\in\cQ}b(P)
  \prod_{P\in\cE_t\setminus\cQ}
  C(p_{i(P)}^\star,P).
  \label{eq:end-factor-expansion}
\end{align}
This identity is the bridge between patch positivity and general
initial laws.
